% fixture: amsart/amsart-pagestyle
% tier: T3
% engine: document
% covers: amsart with \pagestyle{plain}: \maketitle's \topskip 42pt on the first page, which \thispagestyle{empty} after it takes (\ps@empty sets \topskip\normaltopskip again as the page is output); no \@mkboth mark under plain
% issue: #360
\documentclass[11pt]{amsart}
\pagestyle{plain}
\begin{document}
\title{A Title}
\author{Ann Author}
\maketitle
\thispagestyle{empty}
\section{Words}
Typesetting is the composition of text by means of arranging physical type or its digital equivalents. Stored letters and other symbols are retrieved and ordered according to the orthography of a language for visual display. The breaking of paragraphs into lines is a small part of this craft, but it decides how evenly the words are spread and how easily the eye moves from one line to the next.

Typesetting is the composition of text by means of arranging physical type or its digital equivalents. Stored letters and other symbols are retrieved and ordered according to the orthography of a language for visual display. The breaking of paragraphs into lines is a small part of this craft, but it decides how evenly the words are spread and how easily the eye moves from one line to the next.

Typesetting is the composition of text by means of arranging physical type or its digital equivalents. Stored letters and other symbols are retrieved and ordered according to the orthography of a language for visual display. The breaking of paragraphs into lines is a small part of this craft, but it decides how evenly the words are spread and how easily the eye moves from one line to the next.

Typesetting is the composition of text by means of arranging physical type or its digital equivalents. Stored letters and other symbols are retrieved and ordered according to the orthography of a language for visual display. The breaking of paragraphs into lines is a small part of this craft, but it decides how evenly the words are spread and how easily the eye moves from one line to the next.

Typesetting is the composition of text by means of arranging physical type or its digital equivalents. Stored letters and other symbols are retrieved and ordered according to the orthography of a language for visual display. The breaking of paragraphs into lines is a small part of this craft, but it decides how evenly the words are spread and how easily the eye moves from one line to the next.

Typesetting is the composition of text by means of arranging physical type or its digital equivalents. Stored letters and other symbols are retrieved and ordered according to the orthography of a language for visual display. The breaking of paragraphs into lines is a small part of this craft, but it decides how evenly the words are spread and how easily the eye moves from one line to the next.

Typesetting is the composition of text by means of arranging physical type or its digital equivalents. Stored letters and other symbols are retrieved and ordered according to the orthography of a language for visual display. The breaking of paragraphs into lines is a small part of this craft, but it decides how evenly the words are spread and how easily the eye moves from one line to the next.

Typesetting is the composition of text by means of arranging physical type or its digital equivalents. Stored letters and other symbols are retrieved and ordered according to the orthography of a language for visual display. The breaking of paragraphs into lines is a small part of this craft, but it decides how evenly the words are spread and how easily the eye moves from one line to the next.

Typesetting is the composition of text by means of arranging physical type or its digital equivalents. Stored letters and other symbols are retrieved and ordered according to the orthography of a language for visual display. The breaking of paragraphs into lines is a small part of this craft, but it decides how evenly the words are spread and how easily the eye moves from one line to the next.

Typesetting is the composition of text by means of arranging physical type or its digital equivalents. Stored letters and other symbols are retrieved and ordered according to the orthography of a language for visual display. The breaking of paragraphs into lines is a small part of this craft, but it decides how evenly the words are spread and how easily the eye moves from one line to the next.

Typesetting is the composition of text by means of arranging physical type or its digital equivalents. Stored letters and other symbols are retrieved and ordered according to the orthography of a language for visual display. The breaking of paragraphs into lines is a small part of this craft, but it decides how evenly the words are spread and how easily the eye moves from one line to the next.

Typesetting is the composition of text by means of arranging physical type or its digital equivalents. Stored letters and other symbols are retrieved and ordered according to the orthography of a language for visual display. The breaking of paragraphs into lines is a small part of this craft, but it decides how evenly the words are spread and how easily the eye moves from one line to the next.

Typesetting is the composition of text by means of arranging physical type or its digital equivalents. Stored letters and other symbols are retrieved and ordered according to the orthography of a language for visual display. The breaking of paragraphs into lines is a small part of this craft, but it decides how evenly the words are spread and how easily the eye moves from one line to the next.

Typesetting is the composition of text by means of arranging physical type or its digital equivalents. Stored letters and other symbols are retrieved and ordered according to the orthography of a language for visual display. The breaking of paragraphs into lines is a small part of this craft, but it decides how evenly the words are spread and how easily the eye moves from one line to the next.

\end{document}
